\chapter{Zusammenfassen von Mengenfamilien}
\label{b13guide:ch:families}

\section{Ersetzen einzelner Mengen durch ihre Grenzen}


\begin{definition}[Indizierte Infima und Suprema]
\label{b13guide:def:indexed}
Für eine vorgegebene Familie $(x_i)_{i\in I}$ von Elementen von $P$
bedeutet die Indexschreibweise
\[
 \sup_{i\in I}x_i=\sup\{x_i\mid i\in I\},\qquad
 \inf_{i\in I}x_i=\inf\{x_i\mid i\in I\},
\]
sofern die jeweilige Grenze existiert. Wiederholungen von Werten
ändern die zugrunde liegende Menge nicht. Bei leerem $I$ gelten
die Existenzbedingungen für die leere Menge aus
Theorem~\ref{b13guide:thm:empty-extrema}.
\end{definition}

\begin{theorem}[Schranken der Menge der Suprema beziehungsweise Infima]
\label{b13guide:thm:replace-family}
Sei $(S_i)_{i\in I}$ eine Familie von Teilmengen von $P$.
Existiert für jedes $i\in I$ das Supremum $s_i=\sup S_i$, so gilt
\[
 \Hi\Bigl(\bigcup_{i\in I}S_i\Bigr)
       =\Hi(\{s_i\mid i\in I\}).
\]
Existiert für jedes $i\in I$ das Infimum $m_i=\inf S_i$, so gilt
\[
 \Lo\Bigl(\bigcup_{i\in I}S_i\Bigr)
       =\Lo(\{m_i\mid i\in I\}).
\]
Dabei werden nur die genannten inneren Werte vorausgesetzt.
\end{theorem}

\begin{theorem}[Supremum der Suprema; Infimum der Infima]
\label{b13guide:thm:flatten}
Unter der jeweiligen inneren Existenzannahme aus
Theorem~\ref{b13guide:thm:replace-family} gilt
\[
\begin{aligned}
 \sup\Bigl(\bigcup_{i\in I}S_i\Bigr)\text{ existiert}
   &\ \Longleftrightarrow\ \sup\{s_i\mid i\in I\}\text{ existiert},\\
 \inf\Bigl(\bigcup_{i\in I}S_i\Bigr)\text{ existiert}
   &\ \Longleftrightarrow\ \inf\{m_i\mid i\in I\}\text{ existiert}.
\end{aligned}
\]
Im jeweiligen Existenzfall gilt
\[
 \sup\Bigl(\bigcup_{i\in I}S_i\Bigr)=\sup\{s_i\mid i\in I\},\qquad
 \inf\Bigl(\bigcup_{i\in I}S_i\Bigr)=\inf\{m_i\mid i\in I\}.
\]
Auch $I=\varnothing$ ist erlaubt. Dann sind beide zu vergleichenden
Mengen leer; die Existenz des äußeren Wertes wird weiterhin nicht
behauptet. Umgekehrt garantiert die Existenz der Grenze der
Vereinigung allein nicht die Existenz der inneren Grenzen.
\end{theorem}


\section{Mehrfach indizierte Grenzen}

\begin{theorem}[Vertauschen der Indexreihenfolge bei Suprema und Infima]
\label{b13guide:thm:commute-indexed}
Sei $(x_{ij})_{(i,j)\in I\times J}$ eine Familie von Elementen von $P$.
Existieren $\sup_{j\in J}x_{ij}$ für jedes $i\in I$ und
$\sup_{i\in I}x_{ij}$ für jedes $j\in J$, so existiert eines der drei
folgenden äußeren Suprema genau dann, wenn auch die beiden anderen
existieren. Im Existenzfall gilt:
\[
 \sup_{i\in I}\Bigl(\sup_{j\in J}x_{ij}\Bigr)
   =\sup_{(i,j)\in I\times J}x_{ij}
   =\sup_{j\in J}\Bigl(\sup_{i\in I}x_{ij}\Bigr).
\]
Existieren stattdessen sämtliche inneren Infima in beiden
Reihenfolgen, so gilt die entsprechende Aussage für
\[
 \inf_{i\in I}\Bigl(\inf_{j\in J}x_{ij}\Bigr)
   =\inf_{(i,j)\in I\times J}x_{ij}
   =\inf_{j\in J}\Bigl(\inf_{i\in I}x_{ij}\Bigr).
\]
Leere Indexmengen sind erlaubt; die jeweils ausdrücklich verlangten
inneren Existenzbedingungen bleiben auch dann maßgeblich.
\end{theorem}

\begin{theorem}[Gemischte Grenzen: schwache Minimax-Ungleichung]
\label{b13guide:thm:minimax}
Sei $(x_{ij})_{(i,j)\in I\times J}$ eine Familie von Elementen von $P$.
Es mögen alle Werte
\[
 a_i=\inf_{j\in J}x_{ij}\quad(i\in I),\qquad
 b_j=\sup_{i\in I}x_{ij}\quad(j\in J)
\]
sowie $\sup_{i\in I}a_i$ und $\inf_{j\in J}b_j$ existieren.
Dann gilt
\[
 \sup_{i\in I}\Bigl(\inf_{j\in J}x_{ij}\Bigr)
 \ \leq\
 \inf_{j\in J}\Bigl(\sup_{i\in I}x_{ij}\Bigr).
\]
Die Ungleichung kann echt sein; eine Vertauschung mit Gleichheit
folgt aus diesen Voraussetzungen nicht. Auch leere Indexmengen
sind unter den genannten Existenzbedingungen zugelassen.
\end{theorem}

\section{Schranken für gemeinsame Teile}

\begin{theorem}[Infima und Suprema von Durchschnitten vergleichen]
\label{b13guide:thm:intersection-extrema}
Sei $(S_i)_{i\in I}$ eine Familie von Teilmengen von $P$, und
$C=\bigcap_{i\in I}S_i$, mit dem leeren Durchschnitt $C=P$.
Existieren alle $\sup S_i$, ferner $\sup C$ und
$\inf\{\sup S_i\mid i\in I\}$, so gilt
\[
 \sup C\leq\inf\{\sup S_i\mid i\in I\}.
\]
Existieren alle $\inf S_i$, ferner $\inf C$ und
$\sup\{\inf S_i\mid i\in I\}$, so gilt
\[
 \sup\{\inf S_i\mid i\in I\}\leq\inf C.
\]
Gleichheit folgt aus den genannten Voraussetzungen im allgemeinen nicht.
\end{theorem}
